% Separation estimate for the selected finite-volume price test.
\section{A finite-volume separation estimate}
\label{sec:ex-lattice-trace-local-pairing}

We prove the separation estimate that converts a localized probe into a
small two-site correlation. The proof is a three-term bilinear expansion;
all source-specific quantitative inputs are written as precise hypotheses
and audited separately.

Fix a finite periodic gauge configuration space $\Omega$ with normalized
Wilson expectation $\mathbb E_\beta$, and a reflection-positive cell
tiling at selected level $i$ with $n\ge1$ cells. Its admissible
observables are bounded measurable gauge-invariant functions in the
chamber-local domain. For each cell $a$, its copy $C_aF$ is the
translated pullback of $F$ under that cell's reflection; if the reflection
has odd sign parity, the copy is complex-conjugated. Define the root bound
by
\[
 \mathcal R_i(F)=\left|\mathbb E_\beta\prod_{a=1}^n C_aF\right|^{1/n}.
\]
The local positivity estimate supplies the following inequalities for
observables in the chamber-local domain:
\begin{align}
 |\mathbb E_\beta[F(\tau_{a(x)}\omega)G(\tau_{a(y)}\omega)]|
    &\le \mathcal R_i(F)\mathcal R_i(G)\quad(x\ne y),
    \label{eq:ex-lattice-root-product}\\
 \mathcal R_i(F+G)&\le\mathcal R_i(F)+\mathcal R_i(G).
    \label{eq:ex-lattice-root-sum}
\end{align}
Here $x,y$ are coarse sites in the tiling, $a(x),a(y)$ are their fine
anchor sites, and $\omega\in\Omega$ is a gauge configuration. The anchor
map is injective on the two sites under consideration. The first estimate
is the reflected-cell product bound; the second is the root-bound triangle
inequality. We use them only when the observables and all indicated sums
and differences remain in the bounded measurable, gauge-invariant,
chamber-local domain.
The local error scale at selected level $i$ is
\[
 t_i(a)=\exp\!\left(\frac{a}{10^4}(\log p_i)^{48}\right),
 \qquad p_i\longrightarrow\infty.
\]
It satisfies $t_i(a)t_i(c)=t_i(a+c)$ and is increasing in $a$ once $p_i>1$.
The source growth estimate gives $t_i(1)\ge4$ for all sufficiently large
$i$.

\begin{proposition}[Finite-volume two-site error]
\label{prop:ex-lattice-finite-volume-pairing}
Let $P,Q,U,V$ be observables in the chamber-local domain at
selected level $i$, and let $x\ne y$ be coarse sites whose fine anchor
sites are distinct. Suppose
\begin{align}
 \mathcal R_i(P)&\le t_i(5), &\mathcal R_i(Q)&\le t_i(5),
 \label{eq:ex-lattice-pair-main-cost}\\
 \mathcal R_i(Q-V)&\le t_i(-300), &
 \mathcal R_i(P-U)&\le t_i(-200),
 \label{eq:ex-lattice-pair-local-errors}\\
 |\mathbb E_\beta[U(\tau_{a(x)} \omega)V(\tau_{a(y)} \omega)]|&\le t_i(-200).
 \label{eq:ex-lattice-pair-probe}
\end{align}
Then
\[
 |\mathbb E_\beta[P(\tau_{a(x)} \omega)Q(\tau_{a(y)} \omega)]|
 \le t_i(-100)
\]
for all sufficiently large $i$.
\end{proposition}
\begin{proof}
Write $\langle F,G\rangle_{x,y}=\mathbb E_\beta[F(\tau_{a(x)}\omega)G(\tau_{a(y)}\omega)]$. By \eqref{eq:ex-lattice-root-sum} and the assumptions,
\[
 \mathcal R_i(V)\le \mathcal R_i(Q)+\mathcal R_i(V-Q)
       \le t_i(5)+t_i(-300).
\]
The definition gives $\mathcal R_i(-F)=\mathcal R_i(F)$, since negating
$F$ changes the product of its $n$ cell copies only by the overall sign
$(-1)^n$. Hence the assumed bound for $Q-V$ also bounds $V-Q$.
Bilinearity gives the exact identity
\begin{equation}
 \langle P,Q\rangle_{x,y}
 =\langle U,V\rangle_{x,y}
  +\langle P,Q-V\rangle_{x,y}
  +\langle P-U,V\rangle_{x,y}.
 \label{eq:ex-lattice-pair-telescope}
\end{equation}
Indeed, the right hand side equals
$\langle U,V\rangle+\langle P,Q\rangle-\langle P,V\rangle
 +\langle P,V\rangle-\langle U,V\rangle$.
Using \eqref{eq:ex-lattice-root-product} on the last two terms and
\eqref{eq:ex-lattice-pair-main-cost}--\eqref{eq:ex-lattice-pair-probe}
therefore yields
\begin{align*}
 |\langle P,Q\rangle_{x,y}|
 &\le t_i(-200)+t_i(5)t_i(-300)
       +t_i(-200)\bigl(t_i(5)+t_i(-300)\bigr)\\
 &=t_i(-200)+t_i(-295)+t_i(-195)+t_i(-500)\\
 &\le4t_i(-195).
\end{align*}
The last inequality uses monotonicity in the exponent. Since $t_i(1)\ge4$,
\[
 4t_i(-195)\le t_i(1)t_i(-195)=t_i(-194)\le t_i(-100),
\]
again by monotonicity. This proves the conclusion.
\end{proof}

\subsection{Application to a selected response insertion}
\label{subsec:ex-lattice-mileage-application}

For the selected construction, the proposition applies with $P$ equal to
the measured grade-five marker, that is, the one-point observation of the
probe's grade-five label; $Q$ the one-point observation of the selected
grade-four insertion label, which is a separately measured grade-four
field and is not obtained from the marker; $U$ the pullback of the
localized coarse witness supplied by the probe's receipt; and $V$ the
pullback of the bounded coarse approximation supplied, at the same
finite volume and frame level, for that grade-four insertion. The exact source
hypotheses have the following mathematical form: both measured markers have
root bound at most $t_i(5)$; the local approximation satisfies
$\mathcal R_i(Q-V)\le t_i(-300)$; the probe-to-marker difference satisfies
$\mathcal R_i(P-U)\le t_i(-200)$; and the distant-site probe correlation
satisfies $|\langle U,V\rangle_{x,y}|\le t_i(-200)$. The last bound is
the probe's separated-site estimate applied to the coarse approximation
of the insertion; this is legitimate because that approximation also
satisfies the supremum bound $t_i(150000)$, the grade-four window
condition and the room condition required of a test in that estimate.
With these four receipts, the proposition proves the required two-site
price bound $t_i(-100)$.

The proof above discharges the finite-volume assembly estimate, including
its exponent arithmetic. It does not prove production of the four receipts.
In the selected source, their producers are the coefficient-to-root estimate
for the two measured markers, the local approximation and window bound for
the insertion, the probe witness with its distant-site price estimate, and
the finite-volume reflection-positivity seam. The evidence table records
these exact source statements and the remaining producer boundaries. No
continuum or infinite-volume limit is used in this proposition.
